\chapter{Gaussian hierarchies and imported results}\label{sec:inputs}

This section fixes the scalar and matrix objects used in
\cref{part:static,part:dynamics}, and it states, once, the results from the
literature on which the proof of \cref{thm:barrier-730} relies. In
\cref{sec:inp-steps} we define the Gaussian integration step and record its
elementary properties. In \cref{b:sec:scalar} we define the Parisi functional,
the scalar diffusion and the response $\Gamma$, and we state the Parisi formula,
the uniqueness of its minimizer, the scalar estimates and stationarity
identities of OpenAI's barrier preprint \cite{oai-barriers}, and the fact that the support
of the Parisi measure is an interval. Together these imply that $\Gamma(t)=t$ on
$[0,q_\beta]$ (\cref{cor:inp-identity}). \Cref{sec:inp-finite} states the
finite-size estimates of \cite{quantitative-parisi,hexagon-bias}. In
\cref{sec:inp-hierarchies} we define Gaussian hierarchies, their trial values
and the constrained partition functions, we prove the block factorization of
the unperturbed three-replica path (\cref{lem:inp-blocks}), and we state the
matrix comparison and two Gaussian tools of \cite{oai-barriers}.
\Cref{sec:inp-three} states the two- and three-replica pressure bounds of
\cite{oai-barriers}, and \cref{sec:inp-coverage} its energy and rotation
estimates and its bound on the masses of overlap neighborhoods.

We use the scalar objects from \cref{sec:scalar} on. The hierarchies and the
matrix comparison enter \cref{sec:flow,sec:cubic,sec:fractional,sec:static},
the finite-size estimates enter
\cref{sec:fractional,sec:static,sec:references,sec:barrier-proof},
and the bounds on the masses of overlap neighborhoods enter
\cref{sec:references}. We prove
\cref{lem:inp-steps,lem:inp-blocks,cor:inp-identity} and the elementary facts
recorded in the remarks. Every theorem and proposition of this section is
quoted, without proof, from the source named in its heading, or is a consequence
of a proof given there; in the latter case we say so and explain how it is
obtained (\cref{inp:ob-remote}). A note after each result lists the sections
in which it is used.

\paragraph{Imported results and their normalization.}
Apart from standard tools of probability, such as Gaussian integration by
parts, It\^o's formula, and the inequalities of Markov and Chebyshev, the
results stated in this section with a citation are the only mathematical inputs
from the literature in the proof of \cref{thm:barrier-730}. The
inputs used only in \cref{part:fluct} are stated in \cref{sec:fluct,sec:bias}.
The inputs to \cref{part:upper} are listed in \cref{u:sec:scalar}.
The barrier preprint \cite{oai-barriers} uses our normalization
(\cref{tab:model-dictionary}), and it defines constrained pressures as
$n^{-1}\E\log$ of unnormalized constrained sums. Every statement of
\cite{oai-barriers} quoted below holds after the changes of letters listed in
\cref{tab:inp-ob}, apart from the explicitly described consequences of proofs
in \cref{inp:ob-remote}. For each of the other sources, we state, where it is
quoted, only the change of normalization that the quoted result needs;
\cref{sec:model-dictionary} compares all of them with ours.

\begin{table}[ht]
\centering
\small
\caption{Notation of \cite{oai-barriers} and of this paper.}\label{tab:inp-ob}
\begin{tabular}{@{}>{\raggedright\arraybackslash}p{0.25\textwidth}>{\raggedright\arraybackslash}p{0.33\textwidth}>{\raggedright\arraybackslash}p{0.34\textwidth}@{}}
\toprule
Object & In \cite{oai-barriers} & Here\\
\midrule
number of replicas & $d$ & $\ell$\\
constrained pressure & $F_{d,n}(D)$ & $p_{\ell,n}(D)$ (\cref{def:inp-constrained})\\
hierarchy masses & $x_{-1}=0$, $x_k=1$, and $x_{j-1}\le x_j$; the $j$th field is
  integrated at mass $x_{j-1}$ & $0\le m_1\le\dots\le m_N\le1$; the $j$th
  increment is integrated at mass $m_j$ (\cref{def:inp-hierarchy})\\
squared norm, penalty & $|U|^2$, $\int x\,\dd|Q|^2$ & $\norm{U}_F^2$,
  $\int m\,\dd\norm{Q}_F^2$ \eqref{eq:inp-penalty}\\
minimizer & a minimizer $\mu$ & $\mu_\beta$, unique by \cref{inp:unique}\\
scalar diffusion & $X_s$ & $Y_t$\\
scalar quantities & $A(s)$, $d_*(s)$ & $\int_s^1\alpha(t)\dd t$,
  $s\alpha(s)+\int_s^1\alpha(t)\dd t$\\
large-overlap cutoff in \cite[Lemma~4.1 and (4.20)]{oai-barriers} & $h$ & $b_*$\\
scales in \cite[Sections~6--7]{oai-barriers} & $L=n^{1/10000}$, $\rho=L^{-1}$
  & $L_{\rm OB}$, $\rho_{\rm OB}$\\
\bottomrule
\end{tabular}
\end{table}

\paragraph{Conventions.}
Throughout, we fix $\beta>0$. Every statement that involves the
Parisi measure $\mu_\beta$, its support $[0,q_\beta]$, or the response $\Gamma$ of
the minimizer assumes $\beta>1$. Constants $C,c$ and thresholds $n_0$ depend only
on $\beta$ and on the parameters named in the statement where they appear.

\section{Gaussian integration steps}\label{sec:inp-steps}

In this section, $\ell\ge1$ is an integer. We write vectors of $\R^\ell$ in bold,
$\mathbf z\cdot\mathbf y$ and $\abs{\mathbf z}$ for the Euclidean inner product
and norm, and $\mathbf g$ for a standard Gaussian vector in $\R^\ell$. The notation
$\Sigma\succeq0$ means that the symmetric $\ell\times\ell$ matrix $\Sigma$ is
positive semidefinite, and $\Sigma^{1/2}$ denotes its positive semidefinite
square root. A function $f\colon\R^\ell\to\R$ has \emph{linear growth} if it is
measurable and there exists $C<\infty$ such that $\abs{f(\mathbf z)}\le
C(1+\abs{\mathbf z})$ for all $\mathbf z$.

\begin{definition}[Gaussian integration step]\label{def:inp-step}
Let $m\in[0,1]$, let $\Sigma\succeq0$ be an $\ell\times\ell$ matrix, and let $f$
have linear growth. The \emph{Gaussian integration step} at mass $m$ with
covariance $\Sigma$ is
\begin{equation}\label{eq:inp-step}
 T_{m,\Sigma}f(\mathbf z)=
 \begin{cases}
  m^{-1}\log\E\exp\bigl(mf(\mathbf z+\Sigma^{1/2}\mathbf g)\bigr), & m>0,\\
  \E f(\mathbf z+\Sigma^{1/2}\mathbf g), & m=0.
 \end{cases}
\end{equation}
The \emph{tilted Gaussian kernel} is the probability kernel
\begin{equation}\label{eq:inp-kernel}
 \mathcal K_{m,\Sigma,f}(\mathbf z,\dd\mathbf y)
 =\exp\bigl\{m\bigl[f(\mathbf y)-T_{m,\Sigma}f(\mathbf z)\bigr]\bigr\}\,
  N(\mathbf z,\Sigma)(\dd\mathbf y),
\end{equation}
where $N(\mathbf z,\Sigma)$ is the law of $\mathbf z+\Sigma^{1/2}\mathbf g$. For
$m=0$ the exponential factor equals one.
\end{definition}

Note that $T_{m,\Sigma}f$ is the normalization that makes
$\mathcal K_{m,\Sigma,f}(\mathbf z,\cdot)$ a probability measure. A singular
$\Sigma$ is allowed. In that case, $N(\mathbf z,\Sigma)$ is carried by $\mathbf
z+\operatorname{range}\Sigma$, and no inverse of $\Sigma$ is used. A rank-one
covariance $\Sigma=\mathbf v\mathbf v^{\mathsf T}$ integrates a single Gaussian
variable in the direction $\mathbf v$, and $T_{m,0}f=f$.

\begin{lemma}[Elementary properties of the integration step]\label{lem:inp-steps}
Let $m\in[0,1]$, let $\Sigma,\Sigma'\succeq0$ be $\ell\times\ell$ matrices, and
let $f,f'$ have linear growth.
\begin{enumerate}[label=\textup{(\alph*)}]
\item The function $T_{m,\Sigma}f$ has linear growth. If $f\le f'$, then
 $T_{m,\Sigma}f\le T_{m,\Sigma}f'$, and $T_{m,\Sigma}(f+c)=T_{m,\Sigma}f+c$ for
 every $c\in\R$.
\item $T_{m,\Sigma}T_{m,\Sigma'}f=T_{m,\Sigma+\Sigma'}f$.
\item Let $\ell=\ell_1+\ell_2$, write $\mathbf z=(\mathbf z_1,\mathbf z_2)$
 accordingly, and suppose that $f(\mathbf z)=f_1(\mathbf z_1)+f_2(\mathbf z_2)$,
 with $f_1,f_2$ of linear growth,
 and that $\Sigma$ is block diagonal with blocks $\Sigma_1,\Sigma_2$. Then
 $T_{m,\Sigma}f(\mathbf z)=T_{m,\Sigma_1}f_1(\mathbf z_1)+T_{m,\Sigma_2}f_2(\mathbf z_2)$.
\item For every integer $k\ge1$, we have $T_{m/k,\Sigma}(kf)=k\,T_{m,\Sigma}f$.
\item The value $T_{m,\Sigma}f(\mathbf z)$ depends only on the restriction of
 $f$ to $\mathbf z+\operatorname{range}\Sigma$.
\end{enumerate}
\end{lemma}

\begin{proof}
For (a), suppose that $\abs{f(\mathbf z)}\le C(1+\abs{\mathbf z})$, and set
$M=C\norm{\Sigma^{1/2}}\abs{\mathbf g}$, where $\norm{\cdot}$ is the operator
norm. Then $\abs{f(\mathbf z+\Sigma^{1/2}\mathbf g)}\le C(1+\abs{\mathbf
z})+M$. For $m>0$, by Jensen's inequality we have $T_{m,\Sigma}f(\mathbf z)\ge
\E f(\mathbf z+\Sigma^{1/2}\mathbf g)\ge -C(1+\abs{\mathbf z})-\E M$. By H\"older's
inequality, the map $m\mapsto m^{-1}\log\E e^{mM}$ is nondecreasing on $(0,1]$,
which implies that $T_{m,\Sigma}f(\mathbf z)\le C(1+\abs{\mathbf z})+\log\E
e^{M}$. The case $m=0$ is the same with expectations, and the monotonicity and
the translation identity are immediate from \eqref{eq:inp-step}.

For (b) with $m>0$, we have $\exp(mT_{m,\Sigma'}f(\mathbf y))=\E\exp(mf(\mathbf
y+\Sigma'^{1/2}\mathbf g'))$ for an independent standard Gaussian vector $\mathbf
g'$. Applying $T_{m,\Sigma}$, we find
\[
 \exp\bigl(mT_{m,\Sigma}T_{m,\Sigma'}f(\mathbf z)\bigr)
 =\E\exp\bigl(mf(\mathbf z+\Sigma^{1/2}\mathbf g+\Sigma'^{1/2}\mathbf g')\bigr),
\]
and $\Sigma^{1/2}\mathbf g+\Sigma'^{1/2}\mathbf g'$ is centered Gaussian with
covariance $\Sigma+\Sigma'$. The case $m=0$ is the same computation without
exponentials. For (c), if $\Sigma$ is block diagonal, then the two blocks of
$\Sigma^{1/2}\mathbf g$ are independent, and the expectation of
$e^{mf_1}e^{mf_2}$ factorizes (for $m=0$, the expectation of $f_1+f_2$ splits).
Part (d) follows from \eqref{eq:inp-step}, since $(k/m)\log\E e^{(m/k)kf}
=k\,m^{-1}\log\E e^{mf}$ for $m>0$. Part (e) holds because $\mathbf
z+\Sigma^{1/2}\mathbf g\in\mathbf z+\operatorname{range}\Sigma$ almost surely.
\end{proof}
\usedin{\cref{lem:inp-blocks} below, and \cref{sec:flow,sec:cubic,sec:fractional},
in the computations with Gaussian hierarchies.}

\section{The Parisi functional and the scalar diffusion}\label{b:sec:scalar}

A \emph{distribution function} is a right-continuous nondecreasing function
$\alpha\colon[0,1]\to[0,1]$ such that $\alpha(1)=1$. Equivalently,
$\alpha(t)=\mu([0,t])$ for a probability measure $\mu$ on $[0,1]$. We pass
freely between $\alpha$ and $\mu$. We say that $\alpha$ is a \emph{step}
distribution function if there exist $0=t_0<t_1<\dots<t_p=1$ such that $\alpha$
is constant on each interval $[t_{i-1},t_i)$. We call such points
\emph{breakpoints} of $\alpha$, and $m_i=\alpha(t_{i-1})$ denotes the value of
$\alpha$ on $[t_{i-1},t_i)$.

For a step distribution function $\alpha$, we define
$\Phi_\alpha\colon[0,1]\times\R\to\R$ by $\Phi_\alpha(1,z)=\log(2\cosh z)$ and,
for $i=p,p-1,\dots,1$ and $t_{i-1}\le t<t_i$,
\begin{equation}\label{eq:inp-scalar-recursion}
 \Phi_\alpha(t,\cdot)=T_{m_i,\beta^2(t_i-t)}\Phi_\alpha(t_i,\cdot),
\end{equation}
with the scalar integration step of \cref{def:inp-step} ($\ell=1$). On each
interval $(t_{i-1},t_i)$, the function $e^{m_i\Phi_\alpha}$ (or $\Phi_\alpha$
itself, if $m_i=0$) solves the backward heat equation
$-\partial_t\psi=(\beta^2/2)\partial_{zz}\psi$. It follows that $\Phi_\alpha$ is
smooth there and solves the Parisi equation
\begin{equation}\label{b:eq:Parisi-PDE}
 -\partial_t\Phi_\alpha(t,z)=\frac{\beta^2}{2}
 \Bigl(\partial_{zz}\Phi_\alpha(t,z)+\alpha(t)\bigl(\partial_z\Phi_\alpha(t,z)\bigr)^2\Bigr),
 \qquad \Phi_\alpha(1,z)=\log(2\cosh z).
\end{equation}
For a general distribution function $\alpha$, we define $\Phi_\alpha$ as the
uniform limit of $\Phi_{\alpha_k}$ along any sequence of step distribution
functions $\alpha_k$ such that $\int_0^1\abs{\alpha_k-\alpha}\dd t\to0$. By
\cref{inp:ob-scalar} below, the limit exists and does not depend on the
sequence, and \eqref{b:eq:Parisi-PDE} holds in time-integrated form, with the
time derivative existing for almost every $t$. The \emph{Parisi functional} is
\begin{equation}\label{b:eq:Parisi-functional}
 \mathcal P_\beta(\alpha)=\Phi_\alpha(0,0)-\frac{\beta^2}{2}\int_0^1t\alpha(t)\dd t.
\end{equation}
We write $u_\alpha=\partial_z\Phi_\alpha$ and $\chi_\alpha=\partial_{zz}\Phi_\alpha$,
and we define the scalar diffusion and the response by
\begin{equation}\label{b:eq:scalar-diffusion}
 \dd Y_t=\beta\dd B_t+\beta^2\alpha(t)u_\alpha(t,Y_t)\dd t,\qquad Y_0=0,\qquad
 \Gamma_\alpha(t)=\E u_\alpha(t,Y_t)^2,
\end{equation}
where $B$ is a standard Brownian motion. These are the scalar conventions of
\cite[Section~3]{oai-barriers}, where the diffusion is denoted $X$.

\begin{proposition}[{Regularity, stability, and scalar evolution \cite[Lemma~3.1]{oai-barriers}}]\label{inp:ob-scalar}
Let $\beta>0$. There exists $C_\beta<\infty$ such that the following holds for
all distribution functions $\alpha$ and $\widetilde\alpha$.
\begin{enumerate}[label=\textup{(\roman*)}]
\item The limit defining $\Phi_\alpha$ exists and does not depend on the
 approximating sequence. The functions $\Phi_\alpha$, $u_\alpha$ and
 $\chi_\alpha$ are continuous on $[0,1]\times\R$, $\Phi_\alpha(t,\cdot)$ is
 even, and
 \[
  \abs{u_\alpha}\le1,\qquad 0<\chi_\alpha\le1,\qquad
  \abs{\chi_\alpha(t,z)-\chi_\alpha(t,z')}\le C_\beta\abs{z-z'}.
 \]
\item We have
 \[
  \sup_{t,z}\Bigl(\abs{\Phi_\alpha-\Phi_{\widetilde\alpha}}
  +\abs{u_\alpha-u_{\widetilde\alpha}}+\abs{\chi_\alpha-\chi_{\widetilde\alpha}}\Bigr)
  \le C_\beta\int_0^1\abs{\alpha(t)-\widetilde\alpha(t)}\dd t.
 \]
\item The equation \eqref{b:eq:scalar-diffusion} has a unique strong solution. If
 $\alpha$ is a step distribution function with breakpoints $t_0<\dots<t_p$, then
 $(Y_{t_i})_{0\le i\le p}$ is a Markov chain whose transition from $t_{i-1}$ to
 $t_i$ is the tilted Gaussian kernel
 $\mathcal K_{m_i,\beta^2(t_i-t_{i-1}),\Phi_\alpha(t_i,\cdot)}$.
\item For $0\le s\le t\le1$,
 \begin{align*}
  u_\alpha(t,Y_t)-u_\alpha(s,Y_s)&=\beta\int_s^t\chi_\alpha(r,Y_r)\dd B_r,\\
  \E\bigl[\chi_\alpha(t,Y_t)\bigm| Y_s\bigr]&=\chi_\alpha(s,Y_s)
   -\beta^2\int_s^t\alpha(r)\E\bigl[\chi_\alpha(r,Y_r)^2\bigm| Y_s\bigr]\dd r.
 \end{align*}
 Moreover, $\Gamma_\alpha$ is continuously differentiable on $[0,1]$, with
 one-sided derivatives at the endpoints, and
 $\Gamma_\alpha(0)=0$, $\Gamma_\alpha'(t)=\beta^2\E\chi_\alpha(t,Y_t)^2$.
\item Let $0\le a\le1$ and $z\in\R$, and let $Y^{a,z}$ solve
 $\dd Y^{a,z}_t=\beta\dd B_t+\beta^2\alpha(t)u_\alpha(t,Y^{a,z}_t)\dd t$ for
 $a\le t\le1$, with $Y^{a,z}_a=z$. This equation has a unique strong solution,
 and the two identities in \textup{(iv)} hold for $Y^{a,z}$ and $a\le s\le t\le1$.
 If $\alpha$ is a step distribution function, then the transition of $Y^{a,z}$
 across each interval of constancy of $\alpha$ after $a$, with $a$ added to the
 breakpoints, is the tilted Gaussian kernel of \textup{(iii)}. The function
 $\Gamma^{a,z}_\alpha(t)=\E u_\alpha(t,Y^{a,z}_t)^2$ is continuously
 differentiable on $[a,1]$, with $\Gamma^{a,z}_\alpha(a)=u_\alpha(a,z)^2$ and
 $(\Gamma^{a,z}_\alpha)'(t)=\beta^2\E\chi_\alpha(t,Y^{a,z}_t)^2$.
\end{enumerate}
\end{proposition}

We note that, in \cite{oai-barriers}, part (i) collects (3.2) and the evenness of
$\Phi_\alpha$ (so $u_\alpha(t,0)=0$), part (ii) is (3.3), part (iii) is the
statement about (3.4), and part (iv) contains (3.5), (3.7), and (3.8). Taking
expectations in the second identity of (iv), we obtain the unconditional form
\cite[(3.6)]{oai-barriers}. The last sentence of \cite[Lemma~3.1]{oai-barriers}
extends the well-posedness and the identities (3.5)--(3.6) to a diffusion started
at any deterministic time and position. In part (v), the conditional identity
then follows from the Markov property. The statement about the transition
kernels follows from the identification of the kernels on each interval of
constant mass in the proof of \cite[Lemma~3.1]{oai-barriers}, which starts from
an arbitrary position. The statement about $\Gamma^{a,z}_\alpha$ follows from the
martingale identity, It\^o's isometry, and the continuity of $\chi_\alpha$ and
$Y^{a,z}$. The initial value
$\Gamma_\alpha(0)=0$ in (iv) uses the start $Y_0=0$. The constant $C_\beta$
does not depend on $\alpha$.
The proof in \cite{oai-barriers} also bounds $\abs{\partial_z^3\Phi_\alpha}$ by
$2e^{3\beta^2}$ for step distribution functions, and it proves continuity in
$\beta$. We prove bounds on all positive spatial derivatives, for every
distribution function, in \cref{sec:scalar}. For the step approximation and the
spatial regularity of the Parisi evolution, \cite{oai-barriers} credits
Auffinger and Chen \cite[Propositions~1--2]{auffinger-chen-2015}.
\usedin{\cref{sec:scalar}, for all scalar estimates, and
\cref{sec:flow,sec:cubic,sec:fractional,sec:static}, for the stability under a
change of the distribution function and for the passage from step
approximations to the minimizer, and \cref{part:upper}, for the scalar
process and its stochastic identities.}

\begin{theorem}[{Guerra's bound and the Parisi formula \cite[Theorem~3]{guerra-2003}, \cite[Theorems~1--2]{guerra-toninelli-2002},
\cite[Theorem~1.1]{parisi-formula}}]\label{inp:parisi}
Let $\beta>0$. For every $n\ge2$ and every distribution function $\alpha$, we
have $p_n(\beta)\le\mathcal P_\beta(\alpha)$. The limit
$P_\beta=\lim_{n\to\infty}p_n(\beta)$ exists, and
$P_\beta=\inf_\alpha\mathcal P_\beta(\alpha)$, where the infimum is over all
distribution functions. In particular, $\Delta_{n,\beta}=P_\beta-p_n(\beta)\ge0$
for every $n\ge2$.
\end{theorem}

The bound $p_n(\beta)\le\mathcal P_\beta(\alpha)$ is Guerra's
replica-symmetry-breaking bound \cite[Theorem~3]{guerra-2003}, whose proof first
treats step distribution functions and extends to the full class by
\cite[Theorem~1]{guerra-2003}. In our normalization it is
\cite[(2.14)]{oai-barriers}, where it is deduced from \cref{inp:ob-matrix} below
with $\ell=1$ (see \cref{lem:inp-blocks}(a)) and extended to all $\alpha$ by
\cref{inp:ob-scalar}(ii). The existence of the limit was proved by Guerra and
Toninelli \cite[Theorems~1--2]{guerra-toninelli-2002} by an interpolation that
shows superadditivity. The identification of the limit is Talagrand's Parisi
formula \cite[Theorem~1.1, (1.14)]{parisi-formula}, applied at zero external
field with counting measure and covariance function $\xi(q)=\beta^2q^2/2$.
% CHECK: Talagrand (2006) Thm 1.1, eq. (1.14), p. 223; Guerra (2003) Thm 3
% (arXiv cond-mat/0205123v1, p. 7); Guerra--Toninelli (2002) Thms 1--2, eqs. (17)--(18).
\usedin{\cref{sec:model}, for the identification of $P_\beta$ with the
infimum of $\mathcal P_\beta$, and \cref{sec:cubic,sec:fractional,sec:static},
where $P_\beta$ is the baseline of every pressure comparison, and
\cref{part:upper}, for the finite-volume pressure bound.}

\begin{theorem}[{Uniqueness of the Parisi measure \cite[Theorem~2 and Corollary~1]{auffinger-chen-uniqueness-2015}}]\label{inp:unique}
Let $\beta>0$. The functional $\mathcal P_\beta$ has exactly one minimizer among
distribution functions.
\end{theorem}

A minimizer exists by an elementary compactness argument. The probability
measures on $[0,1]$ form a weakly compact set, weak convergence implies $L^1$ convergence of the
distribution functions, and $\mathcal P_\beta$ is continuous in $L^1$ by
\cref{inp:ob-scalar}(ii), as in \cite[proof of Proposition~3.2]{oai-barriers}.
The functional of Auffinger and Chen
\cite[(7)]{auffinger-chen-uniqueness-2015} is $\mathcal P^{\rm AC}_{\xi,h}(\mu)
=\Phi^{\rm AC}_\mu(0,h)$, where $\Phi^{\rm AC}_\mu$ solves the Parisi equation
with coefficient $\xi''$ and terminal value $\log\cosh$. For
$\xi(q)=\beta^2q^2/2$ and $h=0$, we have $\xi''=\beta^2$, and
\cref{rem:inp-logcosh} gives
\[
 \mathcal P_\beta(\alpha)=\log2+\mathcal P^{\rm AC}_{\xi,0}(\mu)
 -\frac{\beta^2}{2}\int_0^1t\alpha(t)\dd t,
\]
which is their full variational objective \cite[(8)]{auffinger-chen-uniqueness-2015}.
They prove that $\mathcal P^{\rm AC}_{\xi,0}$ is strictly convex
\cite[Theorem~2]{auffinger-chen-uniqueness-2015}. Subtracting the linear
functional displayed above preserves strict convexity, and
\cite[Corollary~1]{auffinger-chen-uniqueness-2015} gives uniqueness of the
minimizer. We denote the minimizer by $\alpha_\beta$ and the corresponding
probability measure, the \emph{Parisi measure}, by $\mu_\beta$. By uniqueness,
we may speak of the scalar data of the minimizer; the results of
\cite{oai-barriers} quoted below are stated there for an arbitrary minimizer.
\usedin{every section of \cref{part:static,part:dynamics,part:upper}.}

\paragraph{Standing notation.}
Let $\beta>1$. From now on we write $\alpha=\alpha_\beta$ and $\mu=\mu_\beta$
when no confusion can arise, and $\Phi$, $u$, $\chi$, $Y$, and $\Gamma$ denote
$\Phi_\alpha$, $u_\alpha$, $\chi_\alpha$, the diffusion
\eqref{b:eq:scalar-diffusion} of $\alpha$, and $\Gamma_\alpha$. We also write
$w=\partial_z^3\Phi$; this derivative exists and is bounded by the estimates
of \cref{sec:scalar}. Some sections of \cref{part:static} abbreviate
$b=\beta^2$. In this notation, \eqref{b:eq:Parisi-PDE} and
\eqref{b:eq:scalar-diffusion} read
\begin{equation}\label{b:eq:os-scalar}
\begin{gathered}
 \partial_t\Phi=-\frac b2\bigl(\partial_{zz}\Phi+\alpha(\partial_z\Phi)^2\bigr),\qquad
 \Phi(1,z)=\log(2\cosh z),\\
 \dd Y_t=b\alpha(t)u(t,Y_t)\dd t+\sqrt b\dd B_t,\qquad Y_0=0.
\end{gathered}
\end{equation}

\begin{proposition}[{Support conditions and the mean of $\chi$ \cite[Proposition~3.2]{oai-barriers}}]\label{inp:ob-stationarity}
Let $\beta>0$, and let $\alpha$ be a minimizer of $\mathcal P_\beta$, with
probability measure $\mu$, diffusion $Y$ and response $\Gamma$.
\begin{enumerate}[label=\textup{(\roman*)}]
\item We have $0\in\supp\mu$, and $\supp\mu\subset[0,1-\eps_0]$ for some
 $\eps_0>0$.
\item At every $q\in\supp\mu\cap(0,1)$, we have $\Gamma(q)=q$ and
 $\Gamma'(q)\le1$. Also $\Gamma(0)=0$ and $\Gamma'(0)\le1$.
\item For every $s\in[0,1]$,
 \begin{equation}\label{eq:inp-stationarity-identity}
  \E\bigl[\chi(s,Y_s)+\alpha(s)u(s,Y_s)^2\bigr]=s\alpha(s)+\int_s^1\alpha(t)\dd t,
 \end{equation}
 or equivalently
 \begin{equation}\label{eq:inp-stationarity}
  \E\chi(s,Y_s)=\int_s^1\alpha(t)\dd t+\alpha(s)\bigl(s-\Gamma(s)\bigr).
 \end{equation}
 In particular, $\chi(0,0)=\int_0^1\alpha(t)\dd t$ and
 $\beta\int_0^1\alpha(t)\dd t\le1$. Moreover, $\alpha(s)>0$ for every $s\in(0,1]$.
\end{enumerate}
\end{proposition}

In \cite{oai-barriers}, the right side of \eqref{eq:inp-stationarity-identity}
is written $d_*(s)=1-\int_{(s,1]}t\,\mu(\dd t)$, which equals
$s\alpha(s)+\int_s^1\alpha(t)\dd t$ by Fubini's theorem, and
$\int_s^1\alpha(t)\dd t$ is written $A(s)$. The identities
\eqref{eq:inp-stationarity-identity} and \eqref{eq:inp-stationarity} are
\cite[(3.14)--(3.15)]{oai-barriers}, where they are called stationarity
identities. They hold at every $s$, including atoms of $\mu$. In this paper,
the name \emph{stationarity identities} refers instead to the two identities
\eqref{b:eq:contacts} on $(0,q_\beta)$ described after \cref{cor:inp-identity}. The support conditions in (ii) come from the first variation
\cite[(3.9)]{oai-barriers}. As credited in \cite[Section~3]{oai-barriers}, the
first-variation characterization appears in Chen
\cite[Theorem~2 and Proposition~1]{chen-2017} and in Jagannath and Tobasco
\cite[Proposition~1.1 and Corollary~3.10]{jagannath-tobasco-2017}. Part (ii) does
not describe the support of $\mu$; for $\beta>1$, \cref{inp:full-rsb}
describes it.
\usedin{\cref{cor:inp-identity}, \cref{sec:scalar}, and \cref{sec:static},
where \eqref{eq:inp-stationarity} bounds $\E\chi(q,Y_q)-\int_q^1\alpha(t)\dd t$
at overlaps $q$ with $\abs{q-\Gamma(q)}$ small.}

\begin{theorem}[{Full replica symmetry breaking \cite[Theorem~1.1]{full-rsb}}]\label{inp:full-rsb}
For every $\beta>1$, there exists $q_\beta\in(0,1)$ with the following
properties. We have $\supp\mu_\beta=[0,q_\beta]$. The restriction of $\mu_\beta$
to $[0,q_\beta)$ has a nonnegative density with respect to Lebesgue measure that
is infinitely differentiable on $(0,q_\beta)$ and whose derivatives of every
order extend continuously to $0$. Finally, $\mu_\beta(\{q_\beta\})\in(0,1)$.
\end{theorem}

\cite{full-rsb} works with the Hamiltonian $\beta(2n)^{-1/2}\sum_{i,j=1}^n g_{ij}x_ix_j$,
with $(g_{ij})$ independent standard Gaussians, which equals $\beta H_n$ plus an
independent centered Gaussian of variance $\beta^2/2$ common to all
configurations, that is, $\beta H_n^c$ with $H_n^c$ the completed Hamiltonian of
\cref{sec:model-completion}. The Parisi
functional there is $\log2+u_\mu(0,0)-(\beta^2/2)\int_0^1s\mu([0,s])\dd s$, where
$u_\mu$ solves the Parisi equation with terminal value $\log\cosh x$. By
\cref{rem:inp-logcosh}, this functional equals $\mathcal P_\beta$, and $\mu_\beta$
is the same measure in both papers. We cite version~3 of \cite{full-rsb}. Its
version~4 (October 2026) states the same support conclusion, with a density
that is smooth up to $q_\beta$, and adds a quantitative bound on the atom. We
use only the statement above.
% CHECK: arXiv:2607.11756v3, Theorem 1.1, eqs. (1.4)--(1.5), p. 2; normalization (1.1)--(1.3).
\usedin{\cref{cor:inp-identity}, and through it the stationarity identities of
\cref{sec:scalar} and the choice of the intervals of overlaps in
\cref{sec:cubic,sec:static,sec:references}, and \cref{part:upper}, for the
continuation of the scalar process beyond the support.}

\begin{corollary}[The interval on which $\Gamma(t)=t$]\label{cor:inp-identity}
Let $\beta>1$. Then
\begin{equation}\label{b:eq:identity-interval}
 \Gamma(t)=t\qquad(0\le t\le q_\beta),
\end{equation}
and $\beta^2\E\chi(t,Y_t)^2=1$ for $0<t<q_\beta$. Moreover, $\alpha_\beta(0)=0$,
$\alpha_\beta(t)>0$ for $0<t\le1$, $\alpha_\beta(t)=1$ for $t\ge q_\beta$, and
$\alpha_\beta$ is infinitely differentiable on $[0,q_\beta)$, with one-sided
derivatives at $0$.
\end{corollary}

\begin{proof}
By \cref{inp:full-rsb}, we have $\supp\mu_\beta=[0,q_\beta]$ with $q_\beta<1$.
Then $\supp\mu_\beta\cap(0,1)=(0,q_\beta]$, and \cref{inp:ob-stationarity}(ii)
implies that $\Gamma(q)=q$ for $0<q\le q_\beta$. Also $\Gamma(0)=0$ by
\cref{inp:ob-scalar}(iv), which proves \eqref{b:eq:identity-interval}. Recall
that $\Gamma$ is continuously differentiable with
$\Gamma'(t)=\beta^2\E\chi(t,Y_t)^2$. Differentiating
\eqref{b:eq:identity-interval} on $(0,q_\beta)$, we obtain the second claim. For
$t<q_\beta$, \cref{inp:full-rsb} implies that $\alpha_\beta(t)=\mu_\beta([0,t])$ is
the integral from $0$ to $t$ of a density that is smooth on $[0,q_\beta)$; in
particular $\alpha_\beta(0)=0$. For $t\ge q_\beta$ we have $\alpha_\beta(t)=1$.
Finally, $\alpha_\beta(t)>0$ for $t>0$ because $0\in\supp\mu_\beta$.
\end{proof}

In \cref{sec:scalar}, we prove that the identity $\beta^2\E\chi(t,Y_t)^2=1$ and
its companion $\E w(t,Y_t)^2=2\alpha(t)\E\chi(t,Y_t)^3$ hold on $(0,q_\beta)$
with the regularity needed later. These are the stationarity identities
\eqref{b:eq:contacts}.
\usedin{\cref{sec:scalar}, for the stationarity identities, and
\cref{sec:cubic,sec:static}, for the local comparisons, which take place in
$(0,q_\beta)$, and \cref{part:upper}, for the pre-support scalar identities.}

\begin{remark}[Terminal value $\log\cosh$]\label{rem:inp-logcosh}
Some of the sources quoted here and in \cref{part:fluct}, among them
OpenAI's fluctuation preprint \cite{oai-fluctuations} and \cite{full-rsb}, use
the terminal value $\log\cosh z$ in \eqref{b:eq:Parisi-PDE}. Since
$\log(2\cosh z)=\log2+\log\cosh z$, the two solutions differ by the constant
$\log2$ (by \cref{lem:inp-steps}(a) for step distribution functions, and by
approximation in general). They have the same spatial derivatives, the same
diffusion \eqref{b:eq:scalar-diffusion} and the same response $\Gamma$. Adding
$\log2$ to the functional with terminal value $\log\cosh$ recovers
\eqref{b:eq:Parisi-functional}, and the minimizer does not change. The papers
\cite{quantitative-parisi,hexagon-bias} use exactly
\eqref{b:eq:Parisi-PDE} and \eqref{b:eq:Parisi-functional}.
\end{remark}

\section{Finite-size estimates}\label{sec:inp-finite}

We use three estimates of Aronow and Lopatto on the convergence in the Parisi
formula at a fixed temperature $\beta>1$: an upper bound on the variance of
$F_n$, an upper bound on the mean bias $\Delta_{n,\beta}=P_\beta-p_n(\beta)$,
and a bound on the increase of the Parisi functional when a lower bound is
imposed on the distribution function.

\begin{theorem}[{Upper bound on the variance \cite[Theorem~1.4, (1.11)]{quantitative-parisi}}]\label{inp:qp-rate}
For every fixed $\beta>1$, there exist $C_v<\infty$ and $n_0(\beta)$ such that,
for every $n\ge n_0(\beta)$,
\begin{equation}\label{b:eq:external-rate}
 \Var(F_n)\le C_vn^{7/15}.
\end{equation}
\end{theorem}

The normalization of \cite[(1.1)--(1.2)]{quantitative-parisi} is ours
(\cref{tab:model-dictionary}). The theorem in \cite{quantitative-parisi} also
contains the lower bound $\Var(F_n)\ge c_\beta n^{4/15}$ and the bounds
$c_\beta n^{-11/12}\le P_\beta-p_n(\beta)\le C_\beta n^{-2/3}$ on the mean bias,
which we do not use. The upper bound on the mean bias that we use is
\cref{inp:hb-rate}.
% CHECK: arXiv:2609.09103v1 (v1, 8 Sep 2026), Theorem 1.4, eq. (1.11), p. 5; normalization (1.1), p. 4.
\usedin{\cref{sec:static}, for the lower bound on $Z_n$ in the disorder event
of \cref{b:lem:static-event}; and \cref{sec:bp-strong}, for the same lower
bound in \cref{lem:bp-strong-event}.}

\begin{theorem}[{Upper bound on the mean bias with exponent $7/10$ \cite[Theorem~1.1]{hexagon-bias}}]\label{inp:hb-rate}
For every fixed $\beta>1$, there exists $C_b<\infty$ such that
\[
 0\le P_\beta-p_n(\beta)\le C_b n^{-7/10}\qquad(n\ge2).
\]
\end{theorem}

The normalization of \cite[Section~1.1]{hexagon-bias} is ours, and the lower
bound $0\le P_\beta-p_n(\beta)$ is \cref{inp:parisi}.
% CHECK: hexagon:2610.00177v1, Theorem 1.1, p. 2; model and functional in Section 1.1, (1.1).
\usedin{\cref{sec:static}, for the lower bound on $Z_n$ in the disorder event
of \cref{b:lem:static-event}, together with \cref{inp:qp-rate};
\cref{sec:references}, where it is the input \eqref{eq:ref-mean-input} for
\cref{thm:barrier-730}, and its constant enters the choice of the scale
coefficient in \eqref{b:dyn:R-condition}; and \cref{sec:barrier-proof},
through these choices.}

\begin{proposition}[{The profile $\max\{\alpha,3s\}$ \cite[Lemma~7.2, (7.3), and
Proposition~7.12]{quantitative-parisi}}]\label{inp:qp-floor}
Fix $\beta>1$ and let $\alpha=\alpha_\beta$. There exist constants $T_*>0$,
$0<c_1\le c_2<\infty$, $s_0\in(0,1/3]$, and $C<\infty$, depending only on
$\beta$, such that the following holds.
\begin{enumerate}[label=\textup{(\roman*)}]
\item For $0\le t\le T_*$, we have $c_1t\le\alpha(t)\le c_2t$.
\item For $0<s\le s_0$, the function $\max\{\alpha,3s\}$ is a distribution
 function, and
 \begin{equation}\label{b:eq:external-floor}
  0\le\mathcal P_\beta(\max\{\alpha,3s\})-P_\beta\le Cs^5.
 \end{equation}
\end{enumerate}
\end{proposition}

\begin{remark}[Translation and consequences of \cref{inp:qp-floor}]\label{rem:inp-floor}
In \cite[Section~7.1]{quantitative-parisi}, the scalar problem is rescaled by
$\gamma=\beta\alpha$ and $x_{\rm old}=\beta x$, with rescaled functional
$\mathcal P^{\rm en}_\beta(\gamma)=\beta^{-1}\mathcal P_\beta(\gamma/\beta)$.
By Lemma~7.2 there, applied with temperature threshold $\beta_0=\beta$ to the
rescaled minimizer $\gamma_\beta=\beta\alpha_\beta$, we have
$kq\le\gamma_\beta(q)\le C_0q$ on $[0,T_*]$, which is (i) with $c_1=k/\beta$
and $c_2=C_0/\beta$. For a lower bound $a>0$ on the rescaled profile
(denoted $\lambda$ there), its estimate (7.3) bounds the value of the particular
trial profile $\max\{a,\gamma_\beta\}$ by $\mathcal
P^{\rm en}_\beta(\gamma_\beta)+C_1a^5$ for $0<a\le a_*$, where $a_*>0$ depends
only on $\beta$ (denoted $\lambda_*$ there). With $a=3\beta s$, we have
$\max\{a,\gamma_\beta\}=\beta\max\{3s,\alpha\}$, and we find that
\[
 \mathcal P_\beta(\max\{\alpha,3s\})
 =\beta\,\mathcal P^{\rm en}_\beta(\max\{3\beta s,\gamma_\beta\})
 \le P_\beta+3^5C_1\beta^6s^5\qquad(3\beta s\le a_*),
\]
which is (ii). The same estimate, in the original variables, is
\cite[(7.65)--(7.66)]{quantitative-parisi} with imposed lower bound $3s$. The bound
\eqref{b:eq:external-floor} concerns this particular profile, not only the
minimum over profiles with an atom at zero. For the minimizer of that
constrained problem, Pei \cite[Theorem~1(a) and Proposition~2.3(c)]{pei-fifth-order-2026}
obtains sharp fifth-order asymptotics, which we do not use.

Part (i) localizes the modification. If in addition $3s\le c_1T_*$, then
$\max\{\alpha,3s\}=\alpha$ on $[3s/c_1,1]$, since $\alpha(t)\ge c_1t\ge3s$ for
$3s/c_1\le t\le T_*$ and $\alpha(t)\ge\alpha(T_*)\ge c_1T_*\ge3s$ for $t>T_*$.
On $[0,3s/c_1)$ we have $0\le\max\{\alpha,3s\}-\alpha\le3s$. We conclude
that $\int_0^1\abs{\max\{\alpha,3s\}-\alpha}\dd t\le9s^2/c_1$. The lower bound
$3s$ ensures that the mass $\max\{\alpha,3s\}/3$ of a common block of three
coordinates is at least $s$ (\cref{lem:inp-blocks}(b)), as required in
\cref{sec:fractional}.
\end{remark}
% CHECK: arXiv:2609.09103v1, Lemma 7.2 and (7.3), p. 60; Section 7.1 rescaling, pp. 53--54;
% Proposition 7.12, (7.65)--(7.66), pp. 93--94.
\usedin{\cref{sec:fractional}, for the profile used in the positive fractional
comparison.}

\section{Gaussian hierarchies and constrained pressures}\label{sec:inp-hierarchies}

We now define the finite Gaussian hierarchies that serve as trials in the
pressure comparisons of \cref{part:static}. Their scalar case is the recursion
\eqref{eq:inp-scalar-recursion}. For $\ell\times\ell$ matrices $U,V$ we write
$U\mathbin{:}V=\sum_{a,b}U_{ab}V_{ab}$ and $\norm{U}_F^2=U\mathbin{:}U$.

\begin{definition}[Gaussian hierarchies]\label{def:inp-hierarchy}
Fix $\ell\ge1$.
\begin{enumerate}[label=\textup{(\alph*)}]
\item An \emph{admissible path} of length $N\ge1$ is a pair $(Q,m)$, where
 $Q=(Q_0,\dots,Q_N)$ are symmetric $\ell\times\ell$ matrices with $Q_0=0$ and
 $\Delta Q_j=Q_j-Q_{j-1}\succeq0$ for $1\le j\le N$, and $m=(m_1,\dots,m_N)$
 satisfies $0\le m_1\le m_2\le\dots\le m_N\le1$. We call $Q_N$ the
 \emph{endpoint} of the path and $m_j$ the \emph{mass} of the $j$th increment.
\item For a terminal function $f$ of linear growth, the \emph{partial values} are
 $\Psi_N=f$ and
 \begin{equation}\label{eq:inp-hierarchy-recursion}
  \Psi_{j-1}=T_{m_j,\beta^2\Delta Q_j}\Psi_j\qquad(j=N,N-1,\dots,1),
 \end{equation}
 and we set $\Psi(Q,m;f)=\Psi_0$. For a symmetric $\ell\times\ell$ matrix
 $\lambda$, we define
 \[
  f_\lambda(\mathbf z)=\log\sum_{\epsilon\in\{-1,1\}^\ell}
  \exp\bigl(\mathbf z\cdot\epsilon+\epsilon^{\mathsf T}\lambda\epsilon\bigr),
  \qquad
  \Psi(Q,m,\lambda)=\Psi(Q,m;f_\lambda)(0).
 \]
 In particular, $f_0(\mathbf z)=\sum_{a=1}^\ell\log(2\cosh z_a)$.
\item The \emph{path law} of $(Q,m,\lambda)$ is the law of
 $(\mathbf z_1,\dots,\mathbf z_N,\epsilon)$ defined as follows. Let $\mathbf
 z_0=0$. Given $\mathbf z_0,\dots,\mathbf z_{j-1}$, the vector $\mathbf z_j$ has
 law $\mathcal K_{m_j,\beta^2\Delta Q_j,\Psi_j}(\mathbf z_{j-1},\cdot)$, with
 the partial values of $f=f_\lambda$. Given $\mathbf z_1,\dots,\mathbf z_N$,
 the vector $\epsilon\in\{-1,1\}^\ell$ equals $\epsilon'$ with probability
 $\exp(\mathbf z_N\cdot\epsilon'+\epsilon'^{\mathsf T}\lambda\epsilon'-f_\lambda(\mathbf z_N))$.
 For $0\le l\le N$, the \emph{pair law} sharing the initial segment through level $l$ is
 the law of two such sequences that coincide in $(\mathbf z_1,\dots,\mathbf z_l)$
 and whose continuations, including the terminal draws, are conditionally
 independent given $(\mathbf z_1,\dots,\mathbf z_l)$, each with the
 transitions above. We write $\E_l^{(2)}$ for the corresponding expectation.
\end{enumerate}
\end{definition}

Note that the $j$th increment is integrated at mass $m_j$, that mass zero means
an ordinary Gaussian expectation, and that the terminal Gibbs draw of $\epsilon$
plays the role of a level of mass one. Equivalently, one draws independent increments $\mathbf
z_j-\mathbf z_{j-1}$ with law $N(0,\beta^2\Delta Q_j)$ and tilts the $j$th by
the density $\exp\{m_j[\Psi_j(\mathbf z_j)-\Psi_{j-1}(\mathbf z_{j-1})]\}$
given the past. This is the law of \cite[(2.5)--(2.6)]{oai-barriers}, in which
the continuations of a common initial segment are conditionally independent. By
\cref{lem:inp-steps}(a), every partial value has linear growth, and all
quantities above are finite. The \emph{penalty} of an admissible path is
\begin{equation}\label{eq:inp-penalty}
 \int m\dd\norm{Q}_F^2=\sum_{j=1}^Nm_j\bigl(\norm{Q_j}_F^2-\norm{Q_{j-1}}_F^2\bigr).
\end{equation}
By \cref{lem:inp-steps}(b), merging two consecutive increments with the same
mass, or inserting a zero increment whose mass keeps the masses ordered,
changes neither $\Psi(Q,m,\lambda)$ nor the penalty. \Cref{sec:fractional}
uses a variant of \cref{def:inp-hierarchy} in which all masses are at least a
positive number and the disorder is also integrated at positive mass.

\begin{definition}[Trial values]\label{def:inp-trial}
Let $D$ and $\lambda$ be symmetric $\ell\times\ell$ matrices, and let $(Q,m)$ be
an admissible path. The \emph{trial value} of $(Q,m,\lambda)$ for $D$ is
\begin{equation}\label{eq:inp-trial}
 \mathcal T_D(Q,m,\lambda)=\Psi(Q,m,\lambda)-\lambda\mathbin{:}D
 +\frac{\beta^2}{4}\Bigl(\norm{D-Q_N}_F^2-\int m\dd\norm{Q}_F^2\Bigr).
\end{equation}
If $(Q^{(k)},m^{(k)})$ is a sequence of admissible paths along which the
endpoints, the values $\Psi(Q^{(k)},m^{(k)},\lambda)$, and the penalties
converge, we also call the limit of $\mathcal T_D(Q^{(k)},m^{(k)},\lambda)$ a
trial value for $D$.
\end{definition}

When the endpoint equals $D$ and $\lambda=0$, the trial value is
$\Psi(Q,m,0)-(\beta^2/4)\int m\dd\norm{Q}_F^2$. This is the form of the trial
values in \cref{sec:cubic}. The term $(\beta^2/4)\norm{D-Q_N}_F^2$ is the cost
of an endpoint that differs from the prescribed overlap matrix, and it is used in
\cref{sec:static}.

\begin{definition}[Constrained partition functions]\label{def:inp-constrained}
Let $\ell\ge1$ and $n\ge2$. The \emph{overlap matrix} of
$x^1,\dots,x^\ell\in\Sigma_n$ is $(R(x^a,x^b))_{a,b=1}^\ell$. A symmetric
$\ell\times\ell$ matrix $D$ is \emph{feasible} if it is the overlap matrix of at
least one $\ell$-tuple. For feasible $D$, we define
\[
 Z_{\ell,n}(D)=\sum_{\substack{x^1,\dots,x^\ell\in\Sigma_n\\
 R(x^a,x^b)=D_{ab}\ \forall a,b}}
 \exp\Bigl(\beta\sum_{a=1}^\ell H_n(x^a)\Bigr),\qquad
 p_{\ell,n}(D)=\frac1n\E\log Z_{\ell,n}(D).
\]
\end{definition}

Note that $Z_{\ell,n}(D)/Z_n^\ell$ is the probability that $\ell$ independent samples
from $\pi^J$ have overlap matrix $D$. Since this probability is at most one, we have
$p_{\ell,n}(D)\le\ell p_n(\beta)$. Also $p_{1,n}(1)=p_n(\beta)$. A feasible matrix has unit diagonal
and off-diagonal entries in $\{-1,-1+2/n,\dots,1\}$, and there are at most
$(n+1)^{\ell(\ell-1)/2}$ feasible $\ell\times\ell$ matrices. For two and three
replicas we use the matrices
\begin{equation}\label{eq:inp-triple}
 D_2(q)=\begin{pmatrix}1&q\\q&1\end{pmatrix},\qquad
 D(r,q,\delta)=\begin{pmatrix}1&r+\delta&r-\delta\\r+\delta&1&q\\r-\delta&q&1\end{pmatrix},
\end{equation}
the latter in the coordinate order $(v,x,y)$. If $v,x,y\in\Sigma_n$ have overlap
matrix $D(r,q,\delta)$, then $q=R(x,y)$, $r=(R(v,x)+R(v,y))/2$ and
$\delta=(R(v,x)-R(v,y))/2$.

\begin{remark}[Spin flips and permutations]\label{rem:inp-orientation}
Let $\varpi$ be a permutation of $\{1,\dots,\ell\}$ and let
$\epsilon\in\{-1,1\}^\ell$. For a symmetric $D$, set
$(\mathsf O D)_{ab}=\epsilon_a\epsilon_bD_{\varpi(a)\varpi(b)}$. If $(x^b)_{b\le\ell}$ has
overlap matrix $D$, then $(\epsilon_ax^{\varpi(a)})_{a\le\ell}$ has overlap matrix
$\mathsf OD$,
and this map is a bijection between the two sets of tuples. Since
$H_n(-x)=H_n(x)$, it preserves $\sum_aH_n(x^a)$. We conclude that $\mathsf OD$ is
feasible if and only if $D$ is, and that $Z_{\ell,n}(\mathsf OD)=Z_{\ell,n}(D)$ for every
realization of the disorder. In particular, $Z_{2,n}(D_2(-q))=Z_{2,n}(D_2(q))$,
reversing $v$ gives $Z_{3,n}(D(-r,q,-\delta))=Z_{3,n}(D(r,q,\delta))$, and
exchanging $x$ and $y$ gives $Z_{3,n}(D(r,q,-\delta))=Z_{3,n}(D(r,q,\delta))$.
\end{remark}

The next two results, from \cite[Section~2]{oai-barriers}, are the tools used
to compare constrained pressures with trial values. The first is stated for
a hierarchy with a general terminal function, which is the form used in the
interpolation proofs.

\begin{proposition}[{Differentiation of a hierarchy \cite[Lemma~2.2]{oai-barriers}}]\label{inp:ob-hierarchy-derivative}
Let $\mathcal S$ be a nonempty finite set, let $\omega$ be a random element (the
quenched randomness), let $\mathbf g_1,\dots,\mathbf g_N$ be standard Gaussian
vectors, independent of each other and of $\omega$, and let $0\le m_1\le\dots\le
m_N\le1$. For $s\in\mathcal S$, let
\[
 U(s)=a(s,\omega)+\sum_{i=1}^N\mathbf a_i(s)\cdot\mathbf g_i,
 \qquad\E\max_{s\in\mathcal S}\abs{a(s,\omega)}<\infty,
\]
with deterministic vectors $\mathbf a_i(s)$. Set $\Psi_N=\log\sum_{s\in\mathcal S}e^{U(s)}$, and for $j=N,\dots,1$ set
$\Psi_{j-1}=m_j^{-1}\log\E_je^{m_j\Psi_j}$, or $\Psi_{j-1}=\E_j\Psi_j$ if
$m_j=0$, where $\E_j$ integrates $\mathbf g_j$. Under the path law, given
$(\omega,\mathbf g_1,\dots,\mathbf g_{j-1})$ the law of $\mathbf g_j$ has
density $\exp\{m_j(\Psi_j-\Psi_{j-1})\}$ with respect to its Gaussian law, and
the terminal configuration $S$ has the Gibbs law proportional to $e^{U(S)}$.
For $0\le l\le N$, let $\E_l^{(2)}$ be the expectation under the pair law of
$(S^1,S^2)$ sharing $(\omega,\mathbf g_1,\dots,\mathbf g_l)$. Let $m_0=0$,
$m_{N+1}=1$, and $\Delta m_l=m_{l+1}-m_l$ for $0\le l\le N$. Let $e$ be a fixed vector-valued function on $\mathcal S$, and suppose that
$U$ depends on a parameter $v>0$ in one of the following two ways: either
$\mathbf a_j(s)=\sqrt v\,e(s)$
for some $1\le j\le N$, or $j=0$, $\omega=(\omega',\mathbf g)$ with independent
components, $\mathbf g$ a standard Gaussian vector, and
$a(s,\omega)=a_*(s,\omega')+\sqrt v\,\mathbf g\cdot e(s)$ with
$\E\max_s\abs{a_*(s,\omega')}<\infty$. In both cases $v$ occurs only in the
displayed coefficient. Then, with all other parameters fixed, for $v>0$,
\[
 \frac{\partial}{\partial v}\E\Psi_0
 =\frac12\Bigl(\E\norm{e(S)}^2-\sum_{l=j}^N\Delta m_l\,\E_l^{(2)}e(S^1)\cdot e(S^2)\Bigr).
\]
The identity extends to $v=0$ by integration and continuity.
\end{proposition}

In \cite{oai-barriers}, the quenched root is level $0$ and the $j$th field is
integrated at mass $x_{j-1}$; with $x_{j-1}=m_j$, the weights
$x_l-x_{l-1}$ of \cite{oai-barriers} are our $\Delta m_l$. Under $\E_N^{(2)}$ the two
copies share all fields and only the terminal draws are independent. The
hypotheses on $U$ are those of the applications in \cite{oai-barriers}: the
terminal exponent is affine in the Gaussian variables, and $\mathcal S$ is
finite, so all coefficients are bounded. Gaussian exponential integrability and
the integrability of the root term then make every recursion finite and justify
the differentiation, as in \cite[proof of Lemma~2.2]{oai-barriers}.
\usedin{\cref{sec:fractional}, in the proof of the positive fractional matrix
comparison \cref{b:prop:fractional-matrix}, where the first level has positive
mass.}

\begin{proposition}[{Matrix interpolation bound \cite[Proposition~2.3]{oai-barriers}}]\label{inp:ob-matrix}
Let $n\ge2$ and $\ell\ge1$, let $D$ be a feasible $\ell\times\ell$ matrix, let
$(Q,m)$ be an admissible path of $\ell\times\ell$ matrices, and let $\lambda$ be
a symmetric $\ell\times\ell$ matrix. Then
\begin{equation}\label{eq:inp-matrix}
 p_{\ell,n}(D)\le\mathcal T_D(Q,m,\lambda)
 =\Psi(Q,m,\lambda)-\lambda\mathbin{:}D
 +\frac{\beta^2}{4}\Bigl(\norm{D-Q_N}_F^2-\int m\dd\norm{Q}_F^2\Bigr).
\end{equation}
For fixed $D$ and $\lambda$, the bound passes to every trial value for $D$ in the
sense of \cref{def:inp-trial}.
\end{proposition}

This is \cite[(2.11)]{oai-barriers}, with the constrained pressure
\cite[(2.9)]{oai-barriers} and the path integral \cite[(2.10)]{oai-barriers}. In
\cite{oai-barriers} the path is indexed from $Q_{-1}=0$; inserting the quenched
level $Q_0=0$ allowed there and setting $x_{j-1}=m_j$ identifies the two
conventions, including the penalty. The endpoint $Q_N$ is arbitrary, subject only
to positive semidefinite increments. Its diagonal need not equal
one, and the cost of a mismatch with $D$ is
$(\beta^2/4)\norm{D-Q_N}_F^2$. The proposition belongs to Guerra's interpolation
framework \cite{guerra-2003}. As \cite[Section~2]{oai-barriers} explains,
Talagrand's bounds for coupled replicas
\cite[Propositions~3.4 and~4.4]{talagrand-2007} already retain a cost for an
endpoint that differs from the prescribed overlaps, and Panchenko's vector-spin
formulation \cite{panchenko-2018} proves the asymptotic sharpness of such
bounds. The bound is quenched, and the positive fractional version needed in
\cref{sec:fractional} is proved there.
\usedin{\cref{sec:cubic}, for the trial values of the perturbed three-replica
path, with zero mismatch; \cref{sec:static}, for the comparison at separated
overlaps, with mismatch cost; and \cref{sec:fractional}, which adapts it to
positive fractional moments.}

\begin{lemma}[Block factorization]\label{lem:inp-blocks}
Let $\alpha$ be a step distribution function, and let $0=t_0<t_1<\dots<t_p=1$
contain the points where $\alpha$ jumps. Write $m_i=\alpha(t_{i-1})$ and
$\delta_i=t_i-t_{i-1}$.
\begin{enumerate}[label=\textup{(\alph*)}]
\item \textup{(One coordinate.)} The path with $\ell=1$, $Q_i=t_i$ and masses
 $m_i$ ($1\le i\le p$) is admissible. Its partial values for $f_0$ are
 $\Psi_i=\Phi_\alpha(t_i,\cdot)$, its trial value for $D=1$ and $\lambda=0$ is
 $\mathcal P_\beta(\alpha)$, and under its path law $(\mathbf z_1,\dots,\mathbf
 z_p)$ has the law of $(Y_{t_1},\dots,Y_{t_p})$, where $Y$ is the diffusion
 \eqref{b:eq:scalar-diffusion} of $\alpha$.
\item \textup{(Three coordinates.)} Let $0\le r\le q\le1$ belong to
 $\{t_0,\dots,t_p\}$. In the coordinate order $(v,x,y)$, let
 $\mathbf e=(1,1,1)$ and let $\mathbf e_v,\mathbf e_x,\mathbf e_y$ be the unit
 vectors. We use the following increments of $Q$, with their masses:
 \begin{itemize}
 \item common increments $\delta_i\mathbf e\mathbf e^{\mathsf T}$ at mass
  $m_i/3$, for $t_i\le r$;
 \item $v$-increments $\delta_i\mathbf e_v\mathbf e_v^{\mathsf T}$ at mass $m_i$,
  and $xy$-increments $\delta_i(\mathbf e_x+\mathbf e_y)(\mathbf e_x+\mathbf
  e_y)^{\mathsf T}$ at mass $m_i/2$, for $r<t_i\le q$;
 \item independent increments $\delta_i\mathrm{Id}$ at mass $m_i$, for $t_i>q$.
 \end{itemize}
 Let $(Q,m)$ take the common increments in increasing order of $i$, then the
 $v$- and $xy$-increments in any order that keeps each of these two lists in
 increasing order of $i$ and has nondecreasing masses, and then the
 independent increments in increasing order of $i$. Then $(Q,m)$ is admissible,
 $Q_N=D(r,q,0)$,
 \[
  \Psi(Q,m,0)=3\Phi_\alpha(0,0),\qquad
  \int m\dd\norm{Q}_F^2=6\int_0^1t\alpha(t)\dd t,
 \]
 and $\mathcal T_{D(r,q,0)}(Q,m,0)=3\mathcal P_\beta(\alpha)$.
\item \textup{(Two coordinates.)} Let $q\in\{t_0,\dots,t_p\}$. The path with
 increments $\delta_i(\mathbf e_1+\mathbf e_2)(\mathbf e_1+\mathbf
 e_2)^{\mathsf T}$ at mass $m_i/2$ for $t_i\le q$, followed by $\delta_iI$ at
 mass $m_i$ for $t_i>q$, is admissible, its endpoint is $D_2(q)$, and its trial
 value for $D_2(q)$ and $\lambda=0$ is $2\mathcal P_\beta(\alpha)$.
\item \textup{(Path law in (b).)} For $t_i\le q$, let $Y^v_{t_i}$ be the
 $v$-coordinate of the cumulative field $\mathbf z$ after the common
 increments and the $v$-increments with index at most $i$, and let $Y^x_{t_i}$
 be its $x$-coordinate after the common increments and the $xy$-increments with
 index at most $i$. Under the path law of $(Q,m,0)$, each of
 $(Y^v_{t_i})_{t_i\le q}$ and $(Y^x_{t_i})_{t_i\le q}$ has the law of
 $(Y_{t_i})_{t_i\le q}$, the two coincide for $t_i\le r$, and conditionally on
 $(Y^v_{t_i})_{t_i\le r}$ the families $(Y^v_{t_i})_{r<t_i\le q}$ and
 $(Y^x_{t_i})_{r<t_i\le q}$ are independent, each with the conditional law of
 $(Y_{t_i})_{r<t_i\le q}$ given $Y_r=Y^v_r$.
\end{enumerate}
\end{lemma}

Part (b) with step approximations of $\alpha_\beta$ is the base path of
\cite[Section~4.2]{oai-barriers} and \cite[Section~7, Step~1]{oai-barriers},
and its value is \cite[(7.7)]{oai-barriers}. Part (d) is
\cite[(4.6)]{oai-barriers}. The comparisons of \cref{part:static} perturb this
path near $r$ (\cref{sec:cubic,sec:static}) or replace $\alpha$ by
$\max\{\alpha,3s\}$ (\cref{sec:fractional}).

\begin{proof}
We prove (a) first and use it for the other parts. In (b), we call the common
increments, the $v$- and $xy$-increments, and the independent increments the
first, the middle and the last stage. The proof of (b) has three steps:
admissibility and the endpoint, the value, and the penalty.

\paragraph{Part (a).} Since $\alpha$ is nondecreasing and $m_p\le1$, the path is
admissible. By \eqref{eq:inp-scalar-recursion} at $t=t_{i-1}$, we have
$\Phi_\alpha(t_{i-1},\cdot)=T_{m_i,\beta^2\delta_i}\Phi_\alpha(t_i,\cdot)$, and
backward induction from $\Psi_p=\log(2\cosh\cdot)=\Phi_\alpha(1,\cdot)$ shows that
$\Psi_i=\Phi_\alpha(t_i,\cdot)$ for every $i$. Note that the penalty is
$\sum_im_i(t_i^2-t_{i-1}^2)=2\int_0^1t\alpha(t)\dd t$ and that the endpoint is
$1=D$. Then the trial value is $\Phi_\alpha(0,0)-(\beta^2/2)\int_0^1t\alpha(t)\dd t
=\mathcal P_\beta(\alpha)$. Finally, the transition of $\mathbf z_i$ is
$\mathcal K_{m_i,\beta^2\delta_i,\Phi_\alpha(t_i,\cdot)}$, and
\cref{inp:ob-scalar}(iii) states that this is the transition of $Y$ from
$t_{i-1}$ to $t_i$.

\paragraph{Admissibility and the endpoint in part (b).} Every increment is positive
semidefinite. We check that the masses are nondecreasing. For $t_i\le r$ we
have $t_{i-1}<r$, so that the common masses satisfy $m_i/3\le\alpha(r)/3$. Since
$r$ and $q$ are breakpoints, we see that the first $v$- and $xy$-masses are
$\alpha(r)$ and $\alpha(r)/2$, that the first independent mass is $\alpha(q)$,
and that every $v$- and $xy$-mass is at most $\alpha(q)$. Two nondecreasing
lists can always be merged into one nondecreasing list that keeps the order of
each, and this shows that an order as in (b) exists. The masses of $(Q,m)$ are
then nondecreasing and at most $m_p\le1$. Summing the increments, we obtain
\[
 Q_N=r\,\mathbf e\mathbf e^{\mathsf T}+(q-r)\bigl(\mathbf e_v\mathbf e_v^{\mathsf T}
 +(\mathbf e_x+\mathbf e_y)(\mathbf e_x+\mathbf e_y)^{\mathsf T}\bigr)+(1-q)\mathrm{Id},
\]
which has diagonal entries $1$, entries $r$ in the positions $(v,x)$ and $(v,y)$,
and entry $q$ in the position $(x,y)$. This is the matrix $D(r,q,0)$.

\paragraph{The value in part (b).} We integrate backward. The independent increments
are block diagonal with three one-dimensional blocks, and $f_0$ is a sum of
functions of single coordinates. By \cref{lem:inp-steps}(c) and part (a), after
the independent increments the partial value is $\sum_{a\in\{v,x,y\}}\Phi_\alpha(q,z_a)$.
Let $V=\{\mathbf z\in\R^3:z_x=z_y\}$. Note that the common, $v$-, and
$xy$-increments have ranges in $V$ and that the path starts at $0\in V$. By
\cref{lem:inp-steps}(e), at points of $V$ these steps depend only on
restrictions to $V$, which we write as
functions of $(z_v,z)$, where $z=z_x=z_y$. On $V$, a $v$-increment adds a scalar
Gaussian of variance $\beta^2\delta_i$ to $z_v$ and leaves $z$ fixed, and an
$xy$-increment adds a scalar Gaussian of variance $\beta^2\delta_i$ to $z$ and
leaves $z_v$ fixed. We claim that, during the middle stage, the restriction to
$V$ of the partial value has the form
\begin{equation}\label{eq:inp-middle-form}
 \Phi_\alpha(t_v,z_v)+2\Phi_\alpha(t_{xy},z),
\end{equation}
where $t_v$ (respectively $t_{xy}$) is the largest $t_i\le q$ whose $v$-increment
(respectively $xy$-increment) has not yet been integrated, and $t_v=r$
(respectively $t_{xy}=r$) if there is none. At the end of the middle stage this is
the restriction of $\sum_{a\in\{v,x,y\}}\Phi_\alpha(q,z_a)$, with $t_v=t_{xy}=q$. If the next step
backward is the $v$-increment with $t_i=t_v$, then \cref{lem:inp-steps}(c) and
\eqref{eq:inp-scalar-recursion} turn \eqref{eq:inp-middle-form} into
$\Phi_\alpha(t_{i-1},z_v)+2\Phi_\alpha(t_{xy},z)$. If it is the $xy$-increment with
$t_i=t_{xy}$, then \cref{lem:inp-steps}(c) and (d), with $k=2$, give
$T_{m_i/2,\beta^2\delta_i}\bigl(2\Phi_\alpha(t_{xy},\cdot)\bigr)=2T_{m_i,\beta^2\delta_i}\Phi_\alpha(t_{xy},\cdot)
=2\Phi_\alpha(t_{i-1},\cdot)$. This proves the claim. We conclude that, at the start
of the middle stage, the restriction to $V$ is
$\Phi_\alpha(r,z_v)+2\Phi_\alpha(r,z)$. The
common increments have range spanned by $\mathbf e$. On this line, where
$z_v=z_x=z_y=z$, the function equals $3\Phi_\alpha(r,z)$, and
\cref{lem:inp-steps}(d), with $k=3$, gives
$T_{m_i/3,\beta^2\delta_i}\bigl(3\Phi_\alpha(t_i,\cdot)\bigr)=3\Phi_\alpha(t_{i-1},\cdot)$ at each
common increment. After all common increments, the value at $0$ is
$3\Phi_\alpha(0,0)$.

\paragraph{The penalty in part (b).} During the first stage we have $Q=t\,\mathbf e\mathbf
e^{\mathsf T}$ and $\norm{Q}_F^2=9t^2$, so that this stage contributes
$\sum_{t_i\le r}(m_i/3)\cdot9(t_i^2-t_{i-1}^2)$. During the middle stage, if
the $v$-clock is $t_v$ and the $xy$-clock is $t_{xy}$, then $\norm{Q}_F^2=t_v^2+4t_{xy}^2+4r^2$.
A $v$-increment at mass $m_i$ contributes $m_i(t_i^2-t_{i-1}^2)$, and an
$xy$-increment at mass $m_i/2$ contributes $(m_i/2)\cdot4(t_i^2-t_{i-1}^2)$.
During the last stage only the diagonal changes, and each increment contributes
$3m_i(t_i^2-t_{i-1}^2)$. Every index $i$ contributes
$3m_i(t_i^2-t_{i-1}^2)$ in total, through one increment in the first and last
stages and through two increments in the middle stage. The penalty is
$3\sum_{i=1}^pm_i(t_i^2-t_{i-1}^2)=6\int_0^1t\alpha(t)\dd t$. Since the endpoint
is $D(r,q,0)$, the trial value is
$3\Phi_\alpha(0,0)-(3\beta^2/2)\int_0^1t\alpha(t)\dd t=3\mathcal P_\beta(\alpha)$.

\paragraph{Part (c).} The proof is the same, with the $v$-coordinate removed and
the common stage replaced by the $xy$-increments on $[0,q]$.

\paragraph{Part (d).} Under the path law, the $j$th increment has density
$\exp\{m_j[\Psi_j(\mathbf z_j)-\Psi_{j-1}(\mathbf z_{j-1})]\}$ with respect to its
Gaussian law. By the computation of the value, we see that for a $v$-increment with
index $i$ the exponent is $m_i[\Phi_\alpha(t_i,Y^v_{t_i})-\Phi_\alpha(t_{i-1},Y^v_{t_{i-1}})]$:
the term $2\Phi_\alpha(t_{xy},z)$ of \eqref{eq:inp-middle-form} is the same before and
after the step and cancels. For an $xy$-increment at mass $m_i/2$ the exponent is
$(m_i/2)\cdot2[\Phi_\alpha(t_i,Y^x_{t_i})-\Phi_\alpha(t_{i-1},Y^x_{t_{i-1}})]$,
and for a common increment at mass $m_i/3$ it is
$(m_i/3)\cdot3[\Phi_\alpha(t_i,z)-\Phi_\alpha(t_{i-1},z')]$ along the diagonal.
In each case the moving coordinate makes the scalar transition
$\mathcal K_{m_i,\beta^2\delta_i,\Phi_\alpha(t_i,\cdot)}$, which depends only on its own
current value, and the increments of different branches are driven by
independent Gaussian variables. This proves (d), in view of part (a).
\end{proof}

When $\alpha$ is not a step function, for instance $\alpha=\alpha_\beta$, we
use \cref{lem:inp-blocks} with step distribution functions $\alpha_k\to\alpha$
in $L^1$ whose breakpoints contain $r$ and $q$. By \cref{inp:ob-scalar}(ii),
$\mathcal P_\beta(\alpha_k)\to\mathcal P_\beta(\alpha)$, and the last sentence
of \cref{inp:ob-matrix} applies to the limit. The path-law statement (d) also
passes to the limit, in the following sense. Let $Y^{(k)}$ and $Y$ be the scalar
diffusions of $\alpha_k$ and $\alpha$, driven by the same Brownian motion, and
let $u_k=u_{\alpha_k}$. Since $\abs{u_k}\le1$ and $u$ is $1$-Lipschitz in space
($0<\chi\le1$), by Gronwall's inequality we have
\[
 \sup_{t\le1}\abs{Y^{(k)}_t-Y_t}
 \le\beta^2e^{\beta^2}\bigl(\norm{\alpha_k-\alpha}_{L^1}+\norm{u_k-u}_\infty\bigr),
\]
which tends to zero by \cref{inp:ob-scalar}(ii). For the two branches, we use the
same Brownian path up to $r$ and independent Brownian continuations after $r$.
The branch laws of (d) for $\alpha_k$, evaluated at any fixed finite set of
common breakpoints, therefore converge to the corresponding laws of two
diffusions of $\alpha$ that coincide up to $r$ and are conditionally independent
after $r$ given their common value at $r$.
\usedin{\cref{sec:cubic,sec:fractional,sec:static}, for the value and the law
of the unperturbed path.}

\begin{proposition}[{Gaussian tools \cite[Lemma~2.1]{oai-barriers}}]\label{inp:ob-concentration}
Let $\mathbf g$ be a standard Gaussian vector in $\R^M$, and let
$f\colon\R^M\to\R$ be $L$-Lipschitz. Then, for $a>0$ and $\theta\in\R$,
\begin{equation}\label{eq:inp-concentration}
 \PP\bigl(\abs{f(\mathbf g)-\E f(\mathbf g)}\ge a\bigr)\le2\exp\Bigl(-\frac{a^2}{2L^2}\Bigr),
 \qquad
 \log\E\exp\bigl\{\theta\bigl(f(\mathbf g)-\E f(\mathbf g)\bigr)\bigr\}\le\frac{\theta^2L^2}{2}.
\end{equation}
If a probability density on $\R^M$ is proportional to $e^{-V}$, where $V$ is twice
differentiable and $\nabla^2V\succeq I$, then $\Var(f)\le\E\norm{\nabla f}^2$
under this density for every continuously differentiable $f$ with finite right
side. In particular, the variance of a linear functional $\mathbf a\cdot\mathbf
X$ of a random vector $\mathbf X$ with this density is at most $\abs{\mathbf a}^2$.
\end{proposition}

The bounds \eqref{eq:inp-concentration} are the Gaussian concentration
inequality of Tsirelson, Ibragimov, and Sudakov \cite{tsirelson-ibragimov-sudakov-1976},
and the variance bound is the case $\nabla^2V\succeq I$ of the Brascamp--Lieb
inequality \cite[Theorem~4.1]{brascamp-lieb-1976}. The proof in
\cite{oai-barriers} covers \eqref{eq:inp-concentration} and the case of linear
$f$, under an upper bound on $\nabla^2V$ that holds in its applications. The
linear case enters this paper only through \cref{inp:ob-energy-tail}.
\usedin{\cref{sec:static}, for the concentration of logarithms of constrained
partition functions, and \cref{sec:references}, in the proof of
\cref{b:lem:coverage-extension}.}

\begin{remark}[Lipschitz constants of logarithms of constrained partition functions]\label{rem:inp-lipschitz}
The Gaussian vector in applications of \eqref{eq:inp-concentration} is
$J=(J_{ij})_{i<j}\in\R^{n(n-1)/2}$. Let $\mathcal S\subset\Sigma_n^\ell$ be
nonempty and $F_{\mathcal S}(J)=\log\sum_{(x^a)\in\mathcal
S}\exp(\beta\sum_{a=1}^\ell H_n^J(x^a))$. Then $\partial F_{\mathcal S}/\partial
J_{ij}=\beta n^{-1/2}\langle\sum_ax^a_ix^a_j\rangle_{\mathcal S}$, where
$\langle\cdot\rangle_{\mathcal S}$ is the Gibbs average on $\mathcal S$. Each
derivative has absolute value at most $\beta\ell/\sqrt n$. This implies that
$F_{\mathcal S}$ is Lipschitz with constant $\beta\ell\sqrt{(n-1)/2}$; this is the bound used in
\cite[proof of Lemma~4.2]{oai-barriers}. It applies to $\log Z_{\ell,n}(D)$ and,
with $\ell=1$ and $\mathcal S=\Sigma_n$, to $F_n$.
\end{remark}

\section{Two- and three-replica pressure bounds}\label{sec:inp-three}

In this section, we assume that $\beta>1$, we write $\alpha=\alpha_\beta$, and we let $Y$
and $\Gamma$ be the diffusion and response of $\alpha_\beta$. The three results below are proved in
\cite[Sections~4 and~7]{oai-barriers} from \cref{inp:ob-matrix}, applied to
perturbations of the paths of \cref{lem:inp-blocks}, and from the scalar
identities of \cref{inp:ob-scalar,inp:ob-stationarity}. We import their pressure
bounds, not the probability bounds that \cite{oai-barriers} derives from them at
its own scales. We convert them into probability bounds at our scales in
\cref{sec:static}.

\begin{proposition}[{Pair pressure bound \cite[Lemma~4.3, (4.4)]{oai-barriers}}]\label{inp:ob-pair-gap}
For every $n\ge2$ and every $q\in[0,1]$ such that $D_2(q)$ is feasible,
\[
 p_{2,n}(D_2(q))\le2P_\beta-\frac12\bigl(\Gamma(q)-q\bigr)^2.
\]
\end{proposition}

By \cref{rem:inp-orientation}, the same bound holds for $q\in[-1,0]$ with
$\Gamma(\abs q)-\abs q$ in place of $\Gamma(q)-q$. The proof in
\cite{oai-barriers} uses the path of \cref{lem:inp-blocks}(c) and the terminal
multiplier $\lambda_{12}=\lambda_{21}=L/2$ with $L=q-\Gamma(q)$. The lemma in
\cite{oai-barriers} also contains a probability bound, which we do not use, for
the overlaps $R$ at which $\abs{\Gamma(\abs{R})-\abs{R}}$ exceeds a fixed
positive tolerance.
\usedin{\cref{sec:static}, for the pair overlaps at which $\abs{q-\Gamma(q)}$ is
not small and for the overlaps near one.}

\begin{proposition}[{Signed fixed-scale pressure bound \cite[(4.20)]{oai-barriers}}]\label{inp:ob-fixed-gap}
For every $b_*\in(0,1)$, there exist $r_0,\delta_0,\nu,c_0>0$, depending only on
$\beta$ and $b_*$, with the following property. For every $n\ge2$ and every
feasible $D(r,q,\delta)$ with
\[
 0\le r\le r_0,\qquad\abs\delta\le\delta_0,\qquad b_*\le q\le1,\qquad
 \abs{q-\Gamma(q)}\le\nu,
\]
we have
\[
 p_{3,n}(D(r,q,\delta))\le3P_\beta-c_0\delta^2.
\]
\end{proposition}

In \cite{oai-barriers} the cutoff $b_*$ is called $h$ and the tolerance $\nu$ is
called $\tau$, and (4.20) is proved in
\cite[Sections~4.2--4.4]{oai-barriers} for nonempty constraints. No condition
$r\ge\abs\delta$ is imposed, and the two small overlaps $r\pm\delta$ may have
opposite signs. Reversing $v$ (\cref{rem:inp-orientation}), we obtain the same
bound for $-r_0\le r\le0$. With the conditions on $q$, the bound holds whenever
$\abs{r+\delta},\abs{r-\delta}\le\min\{r_0,\delta_0\}$, since then $\abs r$ and
$\abs\delta$ are at most this minimum. The perturbation in the proof acts on a
fixed interval $[s_1,s_2]\subset(0,b_*)$, and the gain comes from the strict
inequality $\beta^2(\int_{b_*}^1\alpha(t)\dd t)(\int_0^1\alpha(t)\dd t)<1$,
which follows from \cref{inp:ob-stationarity}(iii). The paper \cite{oai-barriers} converts (4.20)
into the probability bound of its Lemma~4.1, with tolerance $2n^{-1/100}$, using
its Lemma~4.2 and its upper bound $Cn^{-1/24}$ on the mean bias. We use neither of these
lemmas.
% CHECK: OB (4.20), printed p. 30.
\usedin{\cref{sec:static}, for the fixed-scale bound \eqref{b:eq:fixed-pressure-gap},
and through it for the signed overlap constraints of
\cref{sec:references,sec:induction}.}

\begin{proposition}[{Narrow-annulus comparison \cite[Section~7, Steps~1--8, (7.30)--(7.31)]{oai-barriers}}]\label{inp:ob-remote}
Let $t_0>0$ satisfy $3t_0\le1$ and $\Gamma(t)=t$ for $t_0/2<t<3t_0$. Let
$\rho_{\rm OB}=n^{-1/10000}$, $h=\rho_{\rm OB}^8$, $\eps=\rho_{\rm OB}^{20}$,
$\theta=\rho_{\rm OB}^2$, and $W=(2(1-\theta)^2+\theta^2)/h$. There exist $C,c>0$
and $n_0$, depending only on $\beta$ and $t_0$, with the following properties.
\begin{enumerate}[label=\textup{(\roman*)}]
\item For $n\ge n_0$ and every feasible $D(r,q,\delta)$ with $t_0/2<r<3t_0$,
 $\mu_\beta([r-h,r+h])\le\rho_{\rm OB}^4$, $q\ge r+\rho_{\rm OB}$ and
 $\eps\le\abs\delta\le2\eps$,
 \begin{equation}\label{eq:inp-remote}
 \begin{split}
  p_{3,n}(D(r,q,\delta))\le{}&3P_\beta
  +\delta^2\Bigl\{\beta^2W\Bigl(\E\chi(q,Y_q)-\int_q^1\alpha(t)\dd t\Bigr)\\
  &+2(1-\theta)\theta\beta^2\alpha(r)\bigl(\beta^2\E[\chi(q,Y_q)\chi(r,Y_r)]-1\bigr)\\
  &+C(\rho_{\rm OB}^4+\theta^2)\Bigr\}
  +C\bigl(\abs\delta^3h^{-3/2}+\delta^4h^{-2}\bigr).
 \end{split}
 \end{equation}
\item For every $t_0/2<r<3t_0$ and $r\le q\le1$,
 \begin{equation}\label{eq:inp-susceptibility-loss}
  \beta^2\E[\chi(q,Y_q)\chi(r,Y_r)]
  =1-\beta^4\int_r^q\alpha(t)\E[\chi(t,Y_t)^2\chi(r,Y_r)]\dd t
  \le1-c(q-r).
 \end{equation}
\end{enumerate}
\end{proposition}

By \cref{cor:inp-identity}, the hypothesis on $t_0$ holds for every $t_0\in(0,q_\beta/3]$,
and finitely many intervals $(t_0/2,3t_0)$ with such $t_0$ cover any compact
subset of $(0,q_\beta)$. The estimate \eqref{eq:inp-remote} is
\cite[(7.30)]{oai-barriers}, obtained there from the step version
\cite[(7.29)]{oai-barriers} by passing to $\alpha_\beta$ at fixed $n,r,q,\delta$.
It is proved by applying \cref{inp:ob-matrix} with $\lambda=0$ to the path of
\cref{lem:inp-blocks}(b), perturbed as follows. With the notation of \cref{lem:inp-blocks}(d), for step
approximations whose breakpoints contain $r-h$, $r$, $r+h$, and $q$, let
\[
 \mathsf Y_0=\frac{(Y^x_r-Y^x_{r-h})-(Y^x_{r+h}-Y^x_r)}{h},\qquad
 \mathsf Y_1=\frac{Y^v_{r+h}-Y^v_r}{h},\qquad
 \mathsf Y=(1-\theta)\mathsf Y_0+\theta\mathsf Y_1,
\]
as linear combinations of the Gaussian increments \cite[(7.8)]{oai-barriers}.
The $x$ field receives $+\delta$ times, and the $y$ field $-\delta$ times, the
contribution of each increment to $\mathsf Y$. The endpoint
becomes $D(r,q,\delta)+\delta^2W(\mathbf e_x-\mathbf e_y)(\mathbf e_x-\mathbf
e_y)^{\mathsf T}$, and the mismatch cost $\beta^2\delta^4W^2$
\cite[(7.13)]{oai-barriers} is part of the last error term. The coefficient $W$
is denoted $K$ in \cite{oai-barriers}. The identity in
\eqref{eq:inp-susceptibility-loss} follows from \cref{inp:ob-scalar}(iv) and
$\Gamma'(r)=1$, and the inequality is \cite[(7.31)]{oai-barriers}.

\Cref{inp:ob-remote} is a consequence of the proof in
\cite[Section~7]{oai-barriers} rather than a verbatim statement there, and we
record how it is obtained. The setup of \cite[Section~7]{oai-barriers} also
imposes $3t_0<q_*$, with $q_*$ the constant of \cref{inp:ob-coverage}, and,
before Step~1, the restriction $\abs{q-\Gamma(q)}\le\rho_{\rm OB}^{13}$. The
Gaussian calculations through (7.30) use neither restriction. Their geometric
requirements are $0<r-h$ and $r+h<q\le1$, which hold under the hypotheses of
(i) for all sufficiently large $n$, and the pair restriction is needed only to
obtain (7.32). In \cite[Section~7]{oai-barriers}, (7.31) is stated for the
pairs $(r,q)$ of Lemma~7.1, but its proof uses only $t_0/2<r<3t_0$ and
$r\le q\le1$. Indeed, let $M$ be such that the event
$\{\beta\sup_{t\le1}\abs{B_t}+\beta^2\le M\}$ has probability $p_M>0$; on it
$\abs{Y_t}\le M$ for all $t\le1$. With $c_M=\min_{0\le t\le1,\,\abs z\le M}
\chi(t,z)>0$, we have $\E[\chi(t,Y_t)^2\chi(r,Y_r)]\ge p_Mc_M^3$, and
$\alpha(t)\ge\alpha(t_0/2)>0$ for $t\ge r$. We may take
$c=\beta^4\alpha(t_0/2)p_Mc_M^3$. After imposing the additional restriction
$\abs{q-\Gamma(q)}\le\rho_{\rm OB}^{13}$, \cite{oai-barriers} deduces from
\eqref{eq:inp-remote} and \eqref{eq:inp-susceptibility-loss} the deficit
$c\delta^2\rho_{\rm OB}^3$ \cite[(7.32)]{oai-barriers} and the probability bound
of its Lemma~7.1. We use neither.
% CHECK: OB Section 7, (7.29)--(7.31), printed pp. 42--43.
\usedin{\cref{sec:static}, for the comparison at separated overlaps in
\eqref{b:eq:remote-pressure-gap} and \eqref{b:eq:remote-tent-estimate}.}

\begin{remark}[Where the scales enter the proof of \eqref{eq:inp-remote}]\label{rem:inp-remote-scales}
\Cref{sec:static} needs a version of \eqref{eq:inp-remote} with $h$ and $\theta$
fixed, and with the oscillation $\alpha(r+h)-\alpha(r-h)$ in place of
$\rho_{\rm OB}^4$. This version is not stated in \cite{oai-barriers}, and we
derive it in \cref{sec:static} by following Steps~1--7 of
\cite[Section~7]{oai-barriers}. In those steps, the specific scales enter only
at the following places.
\begin{enumerate}[label=\textup{(\alph*)}]
\item Step~1 uses $0<r-h$ and $r+h<q$, so that the triangular function on
 $[r-h,r+h]$ that defines the triangular covariance perturbation vanishes at
 $0$ and at $q$.
\item The errors $O(\rho_{\rm OB}^4)$ in \cite[(7.14), (7.15), (7.20), (7.23)--(7.28)]{oai-barriers}
 come from the oscillation bound $\alpha(r+h)-\alpha(r-h)\le\rho_{\rm OB}^4$
 \cite[(7.6)]{oai-barriers}, which follows from the hypothesis on
 $\mu_\beta([r-h,r+h])$. Each of them is at most a constant times
 $\alpha(r+h)-\alpha(r-h)$. The ramp error is $O(1)$ and is multiplied by
 $\theta^2$.
\item Step~4 uses $\delta^2/h\to0$ only to bound
 $\E[\mathsf Y^4e^{\delta^2\mathsf Y^2}]$ through Gaussian exponential moments, since
 $\mathsf Y$ is a centered Gaussian of variance $\beta^2W$ plus a term bounded by
 a constant depending on $\beta$; a bound $\delta^2/h\le c_\beta$ with $c_\beta>0$
 small, together with $h\le1$, suffices.
\item Once the geometric requirements in (a) and the moment restriction in
 (c) are satisfied, the quantitative lower bound $q-r\ge\rho_{\rm OB}$, the
 restriction $\abs{q-\Gamma(q)}\le\rho_{\rm OB}^{13}$, and the relations
 between the powers of $\rho_{\rm OB}$ that define $h$, $\eps$ and $\theta$ are
 used only in Steps~8 and~9.
\end{enumerate}
At fixed $h$ and $\theta$, the last error term of \eqref{eq:inp-remote} is
absorbed only for small amplitudes, so the fixed-parameter version in
\cref{sec:static} holds under a condition $\abs\delta\le\delta_0$. It must be
derived from Steps~1--7 with the oscillation of $\alpha$ kept explicit; it does
not follow by applying \cref{inp:ob-remote}(i) at fixed $h$ and $\theta$.
\end{remark}

\section{Energy, disorder rotation, and overlap-neighborhood masses}\label{sec:inp-coverage}

In this section, we assume that $\beta>1$ and that $J,\widetilde J$ are independent
standard Gaussian vectors in $\R^{n(n-1)/2}$, indexed by pairs $i<j$. We write $Z_n^J$, $\pi^J$, and $H_n^J$
to indicate the disorder. The three results below, from \cite[Sections~5--6]{oai-barriers},
use the off-diagonal Hamiltonian without the auxiliary diagonal variables.

\begin{proposition}[{Upper energy tail \cite[Lemma~5.1]{oai-barriers}}]\label{inp:ob-energy-tail}
There exist constants $e_1\in(0,\beta/2)$ and $c_E>0$ such that, for all
sufficiently large $n$,
\[
 \PP_J\Bigl(\pi^J\{x\in\Sigma_n:H_n^J(x)>ne_1\}>e^{-c_En}\Bigr)\le e^{-c_En}.
\]
\end{proposition}

The proof in \cite{oai-barriers} combines the strict bound
$n^{-1}\E\pi^J(H_n^J)\le e_0<\beta/2$ for all sufficiently large $n$ of
\cite[Lemma~3.3]{oai-barriers}, which
uses $\beta>1$, with the linear case of the variance bound in
\cref{inp:ob-concentration} for the law of $J$ given a Gibbs sample.
\usedin{\cref{sec:references}, in the proof of \cref{b:lem:coverage-extension}.}

\begin{proposition}[{Disorder rotation \cite[Lemma~5.2 and (5.5)]{oai-barriers}}]\label{inp:ob-rotation}
For $\theta\in(0,\pi/2)$, let $J_\theta=\cos\theta\,J+\sin\theta\,\widetilde J$. For every
$n\ge2$, every $s\in\R$ and every $u>0$,
\[
 \E\exp\bigl\{s\bigl(\log Z_n^{J_\theta}-\log Z_n^J\bigr)\bigr\}
 \le\exp\Bigl\{\frac{s^2\theta^2\beta^2(n-1)}4\Bigr\}
\]
and
\[
 \PP\bigl(\log Z_n^{J_\theta}-\log Z_n^J\ge u\bigr)
 \le\exp\Bigl\{-\frac{u^2}{\theta^2\beta^2(n-1)}\Bigr\}.
\]
If $0<t<1/2$ and $J'=\sqrt{1-t}\,J+\sqrt t\,\widetilde J$, then for every $u>0$,
\[
 \PP\bigl(\log Z_n^{J'}-\log Z_n^J\ge u\bigr)\le\exp\Bigl(-\frac{u^2}{2\beta^2nt}\Bigr).
\]
\end{proposition}

The first two bounds are \cite[(5.3)--(5.4)]{oai-barriers}. The last is
\cite[(5.5)]{oai-barriers}, which follows from the second with
$\theta=\arcsin\sqrt t$, since $\theta^2\le2t$ for $0<t<1/2$. All three bounds
are uniform in $\theta$ and $t$. They also apply to rotations whose
size depends on $n$, and they concern the joint law of $(J,\widetilde J)$.
\usedin{\cref{sec:references}, in the proof of \cref{b:lem:coverage-extension}.}

\begin{proposition}[{Overlap-neighborhood mass bound \cite[Proposition~6.1]{oai-barriers}}]\label{inp:ob-coverage}
Let $L_{\rm OB}=n^{1/10000}$. There exists $q_*\in(0,1)$, depending only on
$\beta$, with the following property. For every fixed $D_1>0$ there exist
$D_3=D_3(\beta,D_1)>0$ and events $\mathcal C_n(D_1)$, measurable with respect to
$J$, such that $\PP_J(\mathcal C_n(D_1))\to1$ and, on $\mathcal C_n(D_1)$, for
every $S\subset\Sigma_n$,
\[
 \pi^J\Bigl(\Bigl\{x\in S:\pi^J\{z\in S:\abs{R(x,z)}\ge q_*\}<e^{-D_3L_{\rm OB}}\Bigr\}\Bigr)
 <e^{-D_1L_{\rm OB}}.
\]
\end{proposition}

Note that both masses in the display are masses under $\pi^J$ itself, not under
$\pi^J$ conditioned on $S$. The proposition, titled ``Stretched-exponential coverage'' in
\cite{oai-barriers}, is stated there at the scale $L_{\rm OB}$ only. Its proof uses this scale only through $L_{\rm
OB}\to\infty$ and $L_{\rm OB}/n\to0$, and $q_*$ is chosen before $D_1$ from the
constant $e_1$ of \cref{inp:ob-energy-tail}. \Cref{b:lem:coverage-extension} in
\cref{sec:references} states and proves the corresponding bound for every
deterministic sequence $L_n\to\infty$ with $L_n=o(n)$, by the argument of
\cite[Section~6]{oai-barriers}, using
\cref{inp:ob-energy-tail,inp:ob-rotation,inp:ob-concentration}.
\usedin{\cref{sec:references}, as the statement extended in
\cref{b:lem:coverage-extension}.}
